% !TeX root = ../main.tex
% !TeX encoding = UTF-8

\chapter*{Appendix F: AI Standards Address Alignment Map}
\addcontentsline{toc}{chapter}{Appendix F: AI Standards Address Alignment Map}
\markboth{Appendix F: AI Standards Address Alignment Map}{Appendix F: AI Standards Address Alignment Map}

This map is a reference for General Counsel and governance teams who need to set the facet-level requirements of each major AI governance instrument against the Governance Classification (A·E·C·R·M) introduced in Chapter 14. The minimum address shown for each instrument is a floor: the position below which the instrument's obligations are not met. Where several instruments apply at once, the binding floor on each facet is the most demanding requirement across all of them. It is not the combined address of any single standard.

Enforceable and voluntary instruments are listed separately. Adopting the OECD AI Principles, the UNESCO Recommendation or the EU AI Pact meets no formal compliance obligation. At approximately A2·E2·C1\textendash 2·R2·M2, voluntary commitments provide no legal protection and no defensible position under regulatory scrutiny. When an incident occurs, they will be adduced as evidence that the organisation knew of a standard it had not met.

\section{Enforceable Instruments}

\begingroup
\small
\begin{longtable}[]{@{}
  >{\raggedright\arraybackslash}p{0.24\textwidth}
  >{\raggedright\arraybackslash}p{0.17\textwidth}
  >{\raggedright\arraybackslash}p{0.09\textwidth}
  >{\raggedright\arraybackslash}p{0.42\textwidth}@{}}
\toprule\noalign{}
\textbf{Standard} & \textbf{Minimum address} & \textbf{Pressure facets} & \textbf{Notes} \\
\midrule\noalign{}
\endhead
\bottomrule\noalign{}
\endlastfoot
\textbf{EU AI Act, high-risk systems} (Arts.\ 9, 12, 13, 14, 17, 19) & \textbf{A4·E4·C4·R4·M4} & All five & The only standard requiring Level 4 on all five facets at once. Obligations for Annex III systems are deferred to 2 December 2027 and for Annex I systems to 2 August 2028. Penalties of up to \texteuro 15m or 3\% of worldwide turnover apply. \\
\textbf{EU AI Act, prohibited practices} (Art.\ 5) & \textbf{A4·E4·C3·R4·M3} & A, R & In force since 2 February 2025. Penalties of up to \texteuro 35m or 7\% of worldwide turnover. Any prohibited system found in the inventory is an A4 and R4 obligation. \\
\textbf{EU AI Act, general-purpose AI} (Arts.\ 53 to 55) & \textbf{A3·E3·C3·R3·M3} & None & Applicable since 2 August 2025. Convergence point. Systemic risk assessment and incident reporting for models above the compute threshold. \\
\textbf{DORA} (ICT risk, resilience, third-party providers) & \textbf{A3·E4·C4·R4·M3} & E, C & Third-party ICT risk (E4) and independently verified controls (C4) are the pressure facets. Applies to EU financial entities and their ICT third-party providers. \\
\textbf{FCA SM\&CR} (named accountability, conduct rules) & \textbf{A4·E3·C3·R3·M3} & A & A4 is the critical requirement: named Senior Manager accountability for each material AI system, documented at board level. Applies to FCA-regulated firms. \\
\textbf{UK and EU GDPR}, Art.\ 22 (automated decision-making) & \textbf{A3·E3·C3·R3·M3} & None & Convergence point. Meaningful human involvement and explainability for automated decisions with legal or similarly significant effects. \\
\textbf{NIS2} (critical infrastructure, supply chain) & \textbf{A3·E4·C4·R4·M3} & E, C & Mirrors DORA on E and C for operators of critical infrastructure and essential services. AI systems in critical infrastructure are in scope. \\
\textbf{ISO/IEC 42001:2023} (AI management system) & \textbf{A3·E4·C4·R3·M4} & E, C & E4 requires identification and monitoring of all AI systems, including third-party systems. C4 requires independently verified controls. Certifiable. \\
\textbf{ISO/IEC 27001:2022} (information security management) & \textbf{A3·E3·C4·R4·M4} & C, M & Control integrity (C4) and continuous improvement (M4) are the pressure facets. Certifiable. AI systems are in scope as information assets. \\
\textbf{NIST AI RMF} (Govern, Map, Measure, Manage) & \textbf{A3·E3·C4·R3·M4} & C, M & United States framework with global adoption. Voluntary in the US and referenced in some EU contexts. Operational control (C4) and governance maturity (M4) are the pressure facets. \\
\textbf{UK AI safety framework} (DSIT, safety by design) & \textbf{A4·E3·C3·R3·M3} & A & A4 is required for contestability: citizens must be able to challenge automated decisions that affect them, and named board-level accountability is the mechanism. \\
\textbf{ICO AI guidance} (UK data protection) & \textbf{A3·E3·C3·R3·M3} & None & Convergence point. Reinforces GDPR Art.\ 22 within the UK data protection regime. \\
\end{longtable}
\endgroup

\section{Voluntary Instruments}

\begingroup
\small
\begin{longtable}[]{@{}
  >{\raggedright\arraybackslash}p{0.24\textwidth}
  >{\raggedright\arraybackslash}p{0.17\textwidth}
  >{\raggedright\arraybackslash}p{0.51\textwidth}@{}}
\toprule\noalign{}
\textbf{Standard} & \textbf{Minimum address} & \textbf{Notes} \\
\midrule\noalign{}
\endhead
\bottomrule\noalign{}
\endlastfoot
\textbf{OECD AI Principles} (2023 update) & \textbf{A2·E2·C2·R2·M2} & No enforcement weight and no legal protection. Adoption signals awareness of a governance standard that has not been met. \\
\textbf{UNESCO Recommendation on the Ethics of AI} (2021) & \textbf{A2·E2·C1·R2·M2} & No enforcement weight. C1 because the Recommendation sets principles only and requires no operational governance structure. \\
\textbf{EU AI Pact} (voluntary commitment) & \textbf{A2·E2·C2·R2·M2} & A voluntary commitment made ahead of the EU AI Act. Signing it meets no obligation under the Act and provides no legal defence. \\
\end{longtable}
\endgroup

\section{Sector-Specific Combined Ceilings}

For a regulated organisation, the sector ceiling is found by applying every standard in the sector's regulatory stack at once and taking the most demanding requirement on each facet. The ceiling is the governance destination. The convergence point, A3·E3·C3·R3·M3, is the floor below which no formal standard is met.

\begingroup
\small
\begin{longtable}[]{@{}
  >{\raggedright\arraybackslash}p{0.17\textwidth}
  >{\raggedright\arraybackslash}p{0.17\textwidth}
  >{\raggedright\arraybackslash}p{0.09\textwidth}
  >{\raggedright\arraybackslash}p{0.49\textwidth}@{}}
\toprule\noalign{}
\textbf{Sector} & \textbf{Combined ceiling} & \textbf{Pressure facets} & \textbf{Primary driver} \\
\midrule\noalign{}
\endhead
\bottomrule\noalign{}
\endlastfoot
\textbf{Financial services} & \textbf{A4·E4·C4·R4·M4} & A, E & SM\&CR drives A4. DORA drives E4 and C4. The EU AI Act high-risk regime drives R4 and M4. \\
\textbf{Insurance (Lloyd's market)} & \textbf{A4·E4·C4·R4·M4} & A, C & SM\&CR drives A4. Lloyd's requirements and EU AI Act Art.\ 17 drive C4 independently. DORA applies where ICT dependencies exist in Lloyd's infrastructure. \\
\textbf{Aviation} & \textbf{A4·E4·C4·R4·M4} & C, R & CAA expectations and the EU AI Act high-risk regime (safety-critical decision support) drive C4 and R4. NIS2 drives E4. \\
\textbf{Healthcare} & \textbf{A4·E3·C4·R4·M3} & A, C & CQC expectations and the EU AI Act high-risk regime for medical devices drive A4. Evidence-based controls (C4) are required by both. E and M typically sit at 3. \\
\textbf{Public sector (UK)} & \textbf{A4·E3·C3·R3·M3} & A & The contestability requirement of the UK AI safety framework drives A4. Other facets typically sit at 3: active governance programmes without full independent verification. \\
\textbf{Legal and professional services} & \textbf{A3·E3·C3·R3·M3} & None & The convergence point is the current minimum. Sector regulation is still developing, and A4 is likely to follow as conduct rules extend to AI-assisted legal and advisory output. \\
\end{longtable}
\endgroup

\section{Four Observations}

\textbf{First.} The EU AI Act high-risk regime is the only standard requiring Level 4 on all five facets at once. Every other enforceable standard requires Level 4 on some facets and accepts Level 3 on others. An organisation within its scope faces the most demanding single-instrument requirement in the regulatory landscape.

\textbf{Second.} A4 is required by four independent instruments: the EU AI Act high-risk regime, the EU AI Act prohibited practices, FCA SM\&CR and the UK AI safety framework. For a regulated organisation in the UK or EU, A4 is the binding minimum whichever other instruments apply. Named individual accountability at board level is a concurrent obligation under several regimes, not a choice.

\textbf{Third.} E4 is required by DORA, ISO/IEC 42001 and the EU AI Act high-risk regime. An organisation that self-assesses at E3, with all formal AI systems identified, but has not surfaced vendor-embedded and shadow AI, is operating at E2 in practice. The gap between E3 self-assessment and E2 confirmed address is the most common finding in Digital Alibi Assessments across sectors.

\textbf{Fourth.} Voluntary frameworks sit at approximately A2·E2·C1\textendash 2·R2·M2. Adopting them does not move an organisation towards the convergence point. Building a governance programme around voluntary frameworks, and believing it to be adequate, is the governance equivalent of noting a risk without managing it.

\emph{This map reflects the standards landscape as at October 2026. Regulatory positions change with guidance, enforcement decisions and legislation, and the EU AI Act timetable has already moved once. It is a governance reference, not legal advice, and should be checked with counsel against the current text of each instrument.}
